\documentclass[../DoAn.tex]{subfiles}
\begin{document}
\section{Kết luận}

ĐATN đã nghiên cứu và xây dựng thành công một chatbot tư vấn Luật Hôn nhân và Gia đình theo kiến trúc Agentic GraphRAG, nhằm khắc phục ba hạn chế chính của các công cụ tư vấn pháp luật hiện hành tại Việt Nam: trích dẫn căn cứ pháp lý thiếu chính xác, thiếu cảnh báo hiệu lực và mâu thuẫn pháp lý, cùng thiếu minh bạch luồng suy luận trước khi trả lời.

Qua khảo sát thực tế ở Chương \ref{chapter:Related_works} và Chương \ref{section:1.1}, các sản phẩm như AITracuuluat, AI Luật, Chatbot của Bộ Tư pháp và các mô hình ngôn ngữ lớn phổ biến như Gemini đều có thể trả lời câu hỏi pháp luật ở mức tổng quát, song mỗi hướng tiếp cận bộc lộ điểm yếu khác nhau. AITracuuluat đôi khi viện dẫn căn cứ không liên quan hoặc đã hết hiệu lực. Chatbot của Bộ Tư pháp và AI Luật có độ chính xác và tính cập nhật tốt hơn, nhưng chưa cung cấp cảnh báo trạng thái hiệu lực văn bản, cảnh báo mâu thuẫn giữa các căn cứ chồng chéo và chưa hiển thị chi tiết luồng xử lý nội bộ. Gemini và các LLM tổng quát thường thiếu trích dẫn điều khoản đầy đủ, khó bám sát thời hạn hiệu lực của văn bản quy phạm pháp luật và không có cơ chế kiểm chứng quá trình truy xuất. So với các sản phẩm trên, chatbot của đồ án bổ sung ba khả năng then chốt: truy xuất căn cứ pháp lý qua đồ thị tri thức và bộ mẫu truy vấn Cypher theo miền pháp lý, giúp giảm ảo giác so với LLM thuần; tự động kiểm tra hiệu lực, mở rộng quan hệ sửa đổi, thay thế, mâu thuẫn và đưa cảnh báo vào metadata câu trả lời; minh bạch hóa luồng xử lý qua agent trace trên Chainlit và trang trực quan hóa đồ thị tri thức đã truy xuất. Trên bộ 200 câu hỏi đánh giá tự động, hệ thống đạt Faithfulness 88,2\%, Context Recall 83,5\%, Legal Citation Accuracy 83,7\%, Legal Validity Accuracy 100\% (Bảng \ref{table:experimental_summary}). Các chỉ số này cho thấy chatbot đồ án cải thiện rõ so với LLM tổng quát về mức độ bám căn cứ và kiểm soát hiệu lực, đồng thời vượt trội so với các công cụ thương mại đã khảo sát ở khả năng cảnh báo pháp lý và hiển thị luồng suy luận. Tuy nhiên, chatbot của Bộ Tư pháp và AI Luật vẫn có lợi thế về phạm vi văn bản và độ ổn định triển khai quy mô lớn; đồ án hiện tập trung một phân hệ luật chuyên sâu nên chưa thể thay thế hoàn toàn các nền tảng tra cứu đa lĩnh vực.

Trong phạm vi đề tài, đồ án đã hoàn thành các mục tiêu chính. Thứ nhất, thiết kế và triển khai luồng Agentic Workflow gồm ConversationOrchestrator, RouterAgent, 20 RetrieverAgent theo miền pháp lý, SynthesizerAgent và hai công cụ phản hồi trực tiếp làm rõ hoặc trả lời ngay, kết hợp quản lý trạng thái phiên, bộ nhớ truy xuất và tái sử dụng ngữ cảnh giữa các lượt hỏi đáp. Thứ hai, xây dựng đồ thị tri thức Luật Hôn nhân và Gia đình 2014 cùng hơn 20 văn bản liên quan, mô hình hóa quan hệ giữa thực thể pháp lý và quan hệ giữa các văn bản, kèm bộ Cypher template theo từng chủ đề con. Thứ ba, thiết kế cơ chế truy xuất căn cứ theo thời hạn hiệu lực, chuẩn hóa ngữ cảnh pháp lý và đưa cảnh báo hiệu lực, mâu thuẫn vào pipeline sinh câu trả lời. Thứ tư, xây dựng giao diện chatbot và trang visualize đồ thị tri thức, cho phép người dùng quan sát từng bước định tuyến, truy xuất và kiểm chứng căn cứ đã lấy về. Thứ năm, lập bộ testcase 600 cặp câu hỏi--câu trả lời từ Thư viện Pháp luật, thiết kế bộ chỉ số đánh giá chuyên biệt bên cạnh RAGAS, và triển khai hệ thống trên Docker và Render.

Bên cạnh kết quả đạt được, quá trình thực hiện cũng bộc lộ một số hạn chế. Legal Citation Accuracy (83,7\%) và Invalid Citation Rate (12,5\%) chưa đạt ngưỡng mục tiêu 85\% và dưới 10\%. Một phần nguyên nhân đến từ thiếu hụt quan hệ dẫn chiếu trong đồ thị, giới hạn hiểu ngữ cảnh của LLM với câu hỏi phức tạp hoặc đa ý, và RouterAgent đôi khi chọn sai retriever. Hệ thống chưa đánh giá chất lượng cho câu hỏi thủ tục hành chính và câu hỏi cần kết hợp luật tố tụng dân sự như đã nêu ở phạm vi đề tài. Cơ chế phát hiện mâu thuẫn pháp lý hiện dựa trên quan hệ được khai báo sẵn trên đồ thị và quy tắc xử lý cố định, chưa có lớp suy luận tự động phát hiện xung đột ngoài tập quan hệ đã mô hình hóa. Cảnh báo hiệu lực và mâu thuẫn gắn theo toàn bộ ngữ cảnh truy xuất, nên có thể xuất hiện cảnh báo về điều luật không được nhắc trực tiếp trong câu trả lời cuối cùng.

Qua quá trình thực hiện, đồ án rút ra một số bài học kinh nghiệm. Kết hợp đồ thị tri thức với Agentic RAG giúp kiểm soát nguồn căn cứ tốt hơn so với LLM thuần, nhưng chất lượng cuối cùng vẫn phụ thuộc mạnh vào độ bao phủ của đồ thị và độ chính xác của tầng định tuyến. Thiết kế retriever theo miền pháp lý và Cypher template giúp truy vấn có cấu trúc, song tạo áp lực bảo trì khi mở rộng chủ đề hoặc gặp câu hỏi ngoài mẫu. Việc tách ConversationOrchestrator khỏi RouterAgent không trạng thái là cần thiết để quản lý hội thoại đa lượt, nhưng cũng làm tăng độ phức tạp kiểm thử end-to-end. Cuối cùng, đánh giá tự động trên testcase thực tế kết hợp chỉ số RAGAS và chỉ số pháp lý chuyên biệt là hướng phù hợp để đo lường chatbot tư vấn luật, nhưng vẫn cần bổ sung kiểm thử hộp đen định kỳ cho các tình huống biên.

\section{Hướng phát triển}
\label{section:future-work}

Về hoàn thiện các chức năng đã xây dựng, hướng phát triển trước hết tập trung vào tầng định tuyến và truy xuất. Hiện RouterAgent phụ thuộc hoàn toàn vào mô hình ngôn ngữ lớn và kỹ thuật prompt mô tả chức năng 20 retriever. Khi câu hỏi đa ý, mơ hồ hoặc thuộc ranh giới giữa hai miền pháp lý, router có thể chọn thiếu, thừa hoặc sai retriever. Do đó cần xây dựng tập dữ liệu gán nhãn từ log hội thoại và bộ testcase, rồi fine-tune hoặc huấn luyện mô hình định tuyến chuyên biệt nhằm tăng độ chính xác chọn retriever, giảm phụ thuộc vào prompt dài và cải thiện độ ổn định chi phí suy luận. Tương tự, mỗi RetrieverAgent hiện dùng LLM để phân loại Cypher template và trích xuất tham số điền vào mẫu; sai sót ở bước này dẫn đến truy vấn thiếu căn cứ hoặc lấy nhầm phạm vi thời gian. Hướng phát triển là fine-tune riêng cho từng miền hai nhiệm vụ chọn đúng mẫu truy vấn và trích xuất đúng tham số cùng mốc thời gian sự kiện pháp lý, với dữ liệu huấn luyện tổng hợp từ các lần chạy đúng trên testcase và có sự tham gia của chuyên gia sửa nhãn cho ca lỗi.

Về kỹ thuật truy xuất trên đồ thị, hệ thống hiện dùng Cypher template thiết kế trước theo dạng câu hỏi. Câu hỏi lạ, diễn đạt khác thường hoặc ngoài phạm vi template khiến retriever trả về rỗng và chatbot không trả lời được dù tri thức có thể đã có trên đồ thị. Cần bổ sung các cách truy xuất linh hoạt hơn, chẳng hạn sinh truy vấn Cypher từ ngôn ngữ tự nhiên có kiểm soát schema, truy xuất lai kết hợp template cho câu hỏi quen thuộc và truy vấn sinh động cho câu hỏi ngoài mẫu, hoặc tìm kiếm ngữ nghĩa trên nút và cạnh đồ thị làm bước gợi ý trước khi thực thi truy vấn. Mục tiêu là giữ ưu điểm an toàn và tái lập được của template, đồng thời tăng độ bao phủ câu hỏi.

Đối với cơ chế cảnh báo hiệu lực và mâu thuẫn, xử lý hiện tại đưa cảnh báo dựa trên toàn bộ ngữ cảnh truy xuất, không chỉ các điều luật mà tác tử tổng hợp thực sự viện dẫn trong câu trả lời. Do đó người dùng có thể thấy cảnh báo về văn bản không xuất hiện trong lời tư vấn, gây khó hiểu nếu không đọc agent trace. Hướng cải thiện là đối chiếu metadata cảnh báo với tập trích dẫn thực tế trong câu trả lời, chỉ hiển thị cảnh báo gắn với căn cứ được dùng trong lời tư vấn, đồng thời tách riêng cảnh báo ở mức ngữ cảnh trong panel kiểm chứng cho người dùng muốn xem toàn bộ tri thức liên quan. Về xử lý mâu thuẫn, cơ chế hiện nay vẫn chủ yếu dựa trên quan hệ mâu thuẫn và quan hệ hiệu lực đã khai báo sẵn trên đồ thị cùng quy tắc render cố định, chưa tự phát hiện mâu thuẫn ngữ nghĩa giữa các điều luật chưa được gán quan hệ xung đột. Cần bổ sung lớp phân tích mâu thuẫn tự động trên cặp căn cứ đã truy xuất, kết hợp ràng buộc từ đồ thị và gán nhãn từ chuyên gia pháp lý, rồi đánh giá lại bằng chỉ số phát hiện xung đột trong ngữ cảnh.

Ngoài ra, các hướng mở rộng khác gồm bổ sung đồ thị tri thức và retriever sang thủ tục hành chính, luật tố tụng dân sự liên quan hôn nhân gia đình; tích hợp cập nhật văn bản mới theo chu kỳ; tối ưu chi phí API LLM; và củng cố bộ đánh giá định kỳ khi mở rộng template hoặc fine-tune agent. Tóm lại, chatbot Agentic GraphRAG của đồ án đã chứng minh tính khả thi của việc kết hợp đồ thị tri thức pháp luật với kiến trúc multi-agent để tư vấn Luật Hôn nhân và Gia đình minh bạch và có căn cứ hơn các công cụ đã khảo sát. Các hướng phát triển trên hướng tới giảm phụ thuộc prompt, tăng độ bao phủ truy xuất, làm cảnh báo pháp lý sát với câu trả lời người dùng nhận được, và tiến tới cơ chế phát hiện mâu thuẫn tự động, từ đó nâng cao độ tin cậy trước khi mở rộng sang phân hệ pháp luật khác.

\end{document}
